\chapter{Genetic Algorithms}

Genetic algorithms are a class of optimization algorithms inspired by the process of natural selection. They are used to find approximate solutions to optimization and search problems. The paralelism between the process of natural selection and the genetic algorithm is described in the following concepts:
\begin{itemize}
    \item \ul{Chromosome}: A representation of a solution to the problem. It is a data structure that encodes all the parameters of the solution.
    
    They are often represented as strings of bits, but they can also be represented as arrays of numbers or other data structures depending on the problem being solved.

    \item \ul{Gene}: A part of a chromosome that represents a specific parameter or feature of the solution.
    \item \ul{Alele}: A specific value of a gene. They can be binary (0 or 1), discrete (a finite set of values), or continuous (a range of values).
\end{itemize}

A genetic algorithm typically consists of the following steps:
\begin{enumerate}
    \item \ul{Initialization}: A population of chromosomes is generated randomly or based on some heuristic. The size of the population will be constant throughout the algorithm.
    \item \ul{Selection}: The fitness of each chromosome is evaluated using a fitness function that measures how well the solution represented by the chromosome solves the problem. Chromosomes with higher fitness are more likely to be selected for reproduction.
    \item \ul{Crossover}: Pairs of selected chromosomes are combined to create offspring. This is done by exchanging parts of the chromosomes between the pairs.
    \item \ul{Mutation}: With a small probability, random changes are made to the genes of the offspring to introduce variation into the population.
    \item \ul{Replacement}: The new generation of chromosomes (offspring) replaces the old generation, and the process repeats from the selection step until a stopping criterion is met:
    \begin{itemize}
        \item A solution with a satisfactory fitness level is found.
        \item A maximum number of generations is reached.
        \item The population has converged (i.e., there is little variation in the chromosomes).
    \end{itemize}
\end{enumerate}

During the following sections, we will explore some of these steps in more detail.

\section{Crossover}
Crossover is a genetic operator used to combine the genetic information of two parent chromosomes to generate a new offspring chromosome. There are several types of crossover methods, including:
\begin{itemize}
    \item \ul{Cycle Crossover (CX)}:
    % // TODO: En el ejemplo no es como dice internet. No hay ciclos
    % https://commons.apache.org/proper/commons-math/javadocs/api-3.6.1/org/apache/commons/math3/genetics/CycleCrossover.html
    \item \ul{Order-Based Crossover (OX2)}:
    Fixed positions are selected from the first parent, their values are copied to the offspring in the same positions,
    and the holes are filled with the corresponding values from the second parent.
    \begin{table}[ht]
        \centering
        \begin{tabular}{r|cccccccccc}
            Parent 1 & \textbf{A} & \textbf{A} & B & \textbf{B} & \textbf{A} & A & C & C & \textbf{A} & A \\
            Parent 2 & C & C & A & B & A & B & B & A & A & A \\ \hline
            Offspring & \textbf{A} & \textbf{A} & A & \textbf{B} & \textbf{A} & B & B & A & \textbf{A} & A
        \end{tabular}
        \caption{\centering Example of Order-Based Crossover (OX2), where the fixed positions are 0-1, 3-4 and 8.}
    \end{table}
    \item \ul{Position-Based Crossover (POS)}:
    Fixed positions are selected from the first parent. The values in those positions are copied at the start of the offspring, and the remaining values are filled with the start of the second parent.
    \begin{table}[ht]
        \centering
        \begin{tabular}{r|cccccccccc}
            Parent 1 & \textbf{A} & \textbf{A} & B & \textbf{B} & \textbf{A} & A & C & C & \textbf{A} & A \\
            Parent 2 & C & C & A & B & A & B & B & A & A & A \\ \hline
            Offspring & \textbf{A} & \textbf{A} & \textbf{B} & \textbf{A} & \textbf{A} & C & C & A & B & A
        \end{tabular}
        \caption{\centering Example of Position-Based Crossover (POS), where the fixed positions are 0-1, 3-4 and 8.}
    \end{table}

    \item \ul{Voting Crossover (VX)}:
    % // TODO: Error en las diapositivas?
    % // TODO: CC del inicio no debería ser AA?
    % // Xq coge siempre del 2º? Hace una copia exacta

    \item \ul{Alternating Positions Crossover (AP)}:
    % // TODO: Si ya está, no avanzo a la siguiente posición?
    % // Qué pasa si ya está el de ambos padres?
    % // Qué pasa si sigo hasta terminar los padres y no he terminado el hijo?
    \item \ul{Maximal Preservative Crossover (MPX)}:
    % // TODO: Qué pasa si el intervalo escogido es muy corto? Por ejemplo solo tiene 2?
    \item \ul{Merge Crossover (MX)}:
    % // TODO: Qué segmentos se han elegido?
    % // TODO: Mismas dudas que AP
    \item \ul{Sequential Constructive Crossover (SCX)}:
    % // TODO: No lo entiendo. No es igual que AP?
\end{itemize}

% // TODO: En la diapositiva habla también de Ein-Punkt crossover

\section{Mutation}
Mutation is a genetic operator used to maintain genetic diversity within a population of chromosomes. It introduces random changes to the genes of offspring chromosomes, which can help prevent premature convergence to local optima. There are several types of mutation types, as:
\begin{itemize}
    \item \ul{Silent Mutation}: A mutation that does not change the genotype of the chromosome. For example, if a gene is represented as a binary string, a silent mutation would be changing a 0 to another 0 or a 1 to another 1.
    \item \ul{Neutral Mutation}: A mutation that changes the genotype of the chromosome but does not affect the phenotype (the observable characteristics of the solution).
    \item \ul{Conditional Mutation}: A mutation that also changes the phenotype, but is still considered a valid solution to the problem.
    \item \ul{Letal Mutation}: A mutation that results in a non-viable solution (e.g., a solution that violates constraints of the problem).
\end{itemize}

Depending on the problem being solved, different mutation methods will result in letal mutations, which should be avoided since they do not contribute to the search for optimal solutions. The choice of mutation method and the mutation rate (the probability of mutation occurring) can have a significant impact on the performance of a genetic algorithm. A high mutation rate can lead to excessive randomness and slow convergence, while a low mutation rate may not introduce enough diversity into the population, leading to premature convergence. There are different methods to perform mutation, such as:
\begin{itemize}
    \item \ul{Deletion}: A segment of the chromosome is removed.
    \item \ul{Insertion}: A segment of the chromosome is selected and inserted into another position in the chromosome.
    \item \ul{Duplication}: A segment of the chromosome is duplicated and inserted into another position in the chromosome.
    \item \ul{Inversion}: A segment of the chromosome is reversed.
\end{itemize}
% // TODO: Inversion sí. Pero el resto cambian la longitud del cromosoma, no?
% // TODO: Se pueden aplicar dos en una? Para cambiar un valor...
% // TODO: Hay algunas que no son válidas

\section{Selection}

Selection is the process of choosing which chromosomes from the current population will be used to create the next generation. The selection process is based on the fitness of each chromosome, which is determined by a fitness function that evaluates how well the solution represented by the chromosome solves the problem. A correct balance between exploration and exploitation is crucial for the success of a genetic algorithm:
\begin{itemize}
    \item \ul{Exploration}: The ability of the algorithm to search through a wide range of solutions in the solution space. This is important to avoid getting stuck in local optima and to find the global optimum.
    \item \ul{Exploitation}: The ability of the algorithm to focus on promising areas of the solution space that have already been explored. This is important to refine and improve solutions that are already good.
\end{itemize}

Two important concepts in selection are:
\begin{itemize}
    \item \ul{Selection Pressure}: The degree to which the selection process favors fitter chromosomes over less fit ones. High selection pressure can lead to faster convergence but may also increase the risk of premature convergence to local optima.
    \item \ul{Diversity}: The variety of different chromosomes in the population. Maintaining diversity is important to ensure that the algorithm explores a wide range of solutions and does not get stuck in local optima.
\end{itemize}

There are several selection methods, including:
\begin{itemize}
    \item \ul{Elitist Selection}: The best (for instance, 10\%) chromosome(s) from the current population are directly copied to the next generation (without being subject to crossover or mutation), ensuring that the best solution is preserved.
    \item \ul{Steady State Selection}: Only a small number of chromosomes (e.g., 30\%) are replaced in each generation, while the rest of the population remains unchanged. This allows for a gradual evolution of the population over time.
    
    It should be noted that this approach is just a variation of elitist selection, where the portion of the ``elite'' is much larger than in the standard elitist selection.
    % // TODO: Igual que el anterior, no?
    \item \ul{Tournament Selection}: Fixed $k$ chromosomes are randomly selected from the population, and the one with the highest fitness is chosen for reproduction. This process is repeated until the desired number of chromosomes is selected.
    % // TODO: Qué es el valor de p de las diapositivas?
    \item \ul{Roulette Wheel Selection}: Each chromosome is assigned a probability of selection proportional to its fitness. If $f_i$ is the fitness of chromosome $i$, then:
    \[
        P(i) = \frac{f_i}{\sum\limits_{j=1}^{N} f_j}
    \]
    where $N$ is the total number of chromosomes in the population.   
    It is seen as a roulette wheel where each chromosome occupies a portion of the wheel proportional to its fitness, and the wheel is spun to select chromosomes for reproduction. The process is repeated until the desired number of chromosomes is selected.
    \item \ul{Rank-Based Selection}: Chromosomes are ranked based on their fitness, and selection probabilities are assigned based on their rank rather than their absolute fitness values. The roulette wheel selection method is then applied to these new probabilities. This can help to maintain diversity in the population and prevent premature convergence.
\end{itemize}

It should however be noted that the selection process may combine different methods, such as using elitist selection to preserve the best solutions while also applying tournament selection to select the remaining chromosomes needed for reproduction and reaching the desired population size for the next generation.